\chapter*{Tóm tắt}
\phantomsection
\addcontentsline{toc}{chapter}{Tóm tắt}

Đề tài xây dựng một hệ thống thương mại điện tử thời trang tích hợp trí tuệ nhân tạo nhằm cải thiện trải nghiệm mua sắm trực tuyến. Bên cạnh các nghiệp vụ cơ bản như xem sản phẩm, quản lý giỏ hàng và thanh toán mô phỏng, hệ thống tập trung vào hai chức năng chính: trợ lý phối đồ AI (AI Stylist) và thử đồ ảo (Virtual Try-On -- VTON).

\medskip

AI Stylist tiếp nhận yêu cầu bằng ngôn ngữ tự nhiên và sử dụng Gemini để phân tích ý định. Dữ liệu sản phẩm và tủ đồ cá nhân được truy hồi từ LanceDB bằng tìm kiếm ngữ nghĩa, sau đó hệ thống sinh các phương án phối đồ có cấu trúc để giao diện hiển thị và hỗ trợ thao tác thêm vào giỏ hàng, thay thế sản phẩm hoặc thử đồ.

\medskip

Đối với VTON, các tác vụ sinh ảnh được điều phối bất đồng bộ qua RabbitMQ nhằm tránh chặn luồng yêu cầu web. Dịch vụ AI hỗ trợ CatVTON chạy cục bộ và FLUX.2 thông qua nền tảng Replicate; trạng thái cùng kết quả xử lý được chuyển về hệ thống nghiệp vụ và cập nhật đến giao diện bằng Server-Sent Events. Thực nghiệm cho thấy CatVTON bảo toàn tốt hơn các đặc điểm của ảnh đầu vào, trong khi FLUX.2 có thời gian xử lý ngắn hơn đáng kể.

\medskip

Hệ thống được tổ chức dưới dạng kho mã nguồn hợp nhất, gồm giao diện Next.js, hệ thống nghiệp vụ thương mại điện tử MedusaJS, dịch vụ AI xây dựng bằng FastAPI và các gói mã dùng chung. Kết quả triển khai chứng minh khả năng kết hợp tìm kiếm ngữ nghĩa, mô hình ngôn ngữ lớn, mô hình sinh ảnh và xử lý bất đồng bộ trong một quy trình mua sắm thời trang có tính cá nhân hóa.

\clearpage
